\subsection{Recruitment and summary statistics}

Figure \ref{fig:study_flow} reports the dates of all surveys, as well as the number of participants taking them, 
the number of participants eligible for and taking follow-up surveys, and the number with 
valid data. 8,003 people completed the baseline survey, of which 4,525 were flu vaccine 
hesitant and invited to the main survey, which 3,651 began, resulting in 3,525 participants 
after exclusions. (Sample sizes by trial arm, and counts of non-missing secondary outcomes, 
are in Appendix Table \ref{tab:counts}.) All participants who completed the main survey were invited to 
the follow-up survey; 3,210 began the survey and after exclusions, 
we end-up with 2,993 completions. Attrition between the main and follow-up survey mean 
that we are missing vaccination data for about 15 percent of participants. 

Table \ref{tab:balance} reports summary statistics for the analysis sample, by treatment arm.  
The sample is well-balanced in the sense that the differences across arm are small and 
typically insignificant for individual covariates (and jointly insignificant). 
The table also indicates several important features of the sample. First the typical 
participant—selected for flu vaccine hesitancy—believes that a side effect is at least likely 
if they vaccinate, and two thirds do not intend to vaccinate at baeline (the remainder say 
they may or may not vaccinate). Despite their hesitancy, a majority have previously vaccinated 
against flu or COVID, and a substantial number indicate vaccinating and having a severe 
reaction. Somewhat more than half follow their doctor’s advice on vaccines, and fewer than 
half trust the government’s recommendations on vaccinations.  About one in six reports having 
a health condition which would make them vulnerable to flu. Like other samples
drawn from online panels, our sample is younger, better educated, and higher income
than the typical American \citep{stantcheva2023run}.

\subsection{Predictors of personal side effect beliefs}

Side effect beliefs are highly pessimistic compared to clinical trial outcomes.
In our control group, the average participant believes 20 percent of vaccine-arm trial 
participants experienced a severe advsere event (against 3 percent in the placebo arm); 
in fact roughly 1.7 percent did. Likewise the average control group participant believes 
vaccinating would increase her personal side effect risk by 15 percentage points. These high 
rates are not a consequence of a general tendency to overestimate low probabilities;
they are determined by the {\em differnece} in probabilities between vaccinated and 
unvaccinated states. Further, \citet{sacks2025preferences} show that vaccine confident
people estimate mean personal side effect risk to be near zeor. 

Despite this pessimistic bias, in our sampel of vaccine hesitant participants, 
short run personal side effect beliefs are sensibly 
correlated with other beliefs and experiences. To show this, we regress personal side effect b
eliefs on other variables, individually and jointly; we report the results in Table 
\ref{tab:delta_predictors}

The first column shows a strong but imperfect correlation between beliefs about 
personal side effects and beliefs about trial outcomes. Each 1 percentage point increase
in belief about side effects in the trial is associated with an additional 0.6 percentage 
point increase in beleived personal risk. People who think side effects are likely in general
think they are likely to experience side effects themselves. But people also recognize
that their own experience may differ from the trial one. 

The next columns show the association between personalside effect beliefs
and vaccine experience as well as trust in government and following doctors advice. 
Relative to people who never got got a flu vacicne or COVID vaccine, 
people who experienced no vaccine reaction have lower persnal beliefs. People experiencing 
a severe flu vacicne reaction predict they will be more likely to have severe reactions 
to the next vaccine. There is little extrapolation from COVID vaccine
to flu vaccine conditional on flu vaccine experience.
Mistrust of government vaccine advice and uwillingness to follow a doctor's advice on
vaccines are strongly associated with more pessimistic beliefs, as we would expect for 
two variables which are signals of vaccine hesitancy.%
\footnote{Versions of these items are included on commonly used scales to assess 
vaccine hesitancy \citep{larson2015measuring}.}
The final column, which includes all predictors, shows that personal experiences load onto 
personal risk rather than general risk because the coefficients are similar when we condition 
on beliefs about trial outcomes, especially severe prior reactions. 

\subsection{Main experiment results show scientific information can reduce perceived risks}

We present our main, pre-specified treatment effect estimates in Table \ref{tab:treat_effects}. 
Column (1) shows that all treatment arms provided salient information. 
In our control group, the average belief about the vaccine arm severe adverse event rate 
is 20 percentage points. The industry arm reduced this belief by 8 percentage points, 
and the academic and representative arms reduced it by 4-5 percentage points. 

The next column shows effects for our main outcome, personal side effect beliefs. 
The industry and academic arms reduce this belief by about 3 percentage points 
on a base of 15 percentage points in the control group. We do not find a significant 
change in beliefs in the representative arm. The representative intervention appeared 
unsuccessful in the sense that our intention was to raise perceived personal relevance, 
but in fact  participants found this arm less relevant 
(Appendix Table \ref*{tab:trust_relevance}).

These estimates suggest relatively low learning rates. The learning rate is 
the effect of a 1 percentage point information shock on belief updating 
\citep{haaland2023designing}. In our context, the information shock is that the 
side effect rate is roughly 0, the control belief is 15 percentage points, and the 
belief updating in industry and academic arm is about 3 percenatage points, so the learning 
rate is 20 percent.
This is on the lower end of the range of prior estimated learning rates 
\citep{haaland2023designing}. One reason for the low learning rate is that people 
do not fully extrapolate from the trial experience 
to their personal risk; personal risk updating as a share of the trial outcome updating varies 
from roughly 0 (representative arm) to 41\% (industry arm) to 70\% (academic arm). 
Another reason may be that our information was quantitatively vague.
Overall, therefore, scientific information provision can shift beliefs about personal side 
effect risks among the vaccine hesitant, albeit incompletely.

The third column shows how the change in beliefs translated into intentions to vaccinate. 
The academic treatment raises vaccination intentions by 3 percentage points, suggesting that 
reducing side effect concerns can raise vaccine intentions. However, the industry treatment arm 
does not increase vaccine intentions despite reducing side effect perceptions, and the personal 
arm increased intentions despite not increasing the personal side effect belief. We show below 
that this surprising pattern of effects on intentions is explained in part by heterogeneous 
responses by trust in the provided information.

The final two columns show effects on vaccine behavior. Column (4) reports effects on clicking 
a link to schedule a vaccination. The point estimates are small and insignificant, and overall 
few people clicked on the link. The last column shows effects on self-reported vaccination. 
The point estimates range from -1.4 percentage points for the industry arm to about 
2 percentage points for the academic and personal arms. The estimates are imprecise; 
the 95\% confidence intervals include both zero and the effect size for intentions. 
Thus our results are uninformative for vaccine behaviors. 

\subsection{Heterogeneous effects by prior belief and views of university research point to 
belief updating}

Models of belief updating predict larger effects for participants with more biased priors
and for participants who find the information source more credible. Heterogeneity analyses 
confirm both these predictions in our data. We caution that we did not 
prespecify specific heterogeneity analyses. 

To show larger updating for participants with worse priors, we estimate separate regressions
with subsamples defined by prior side effect beliefs.  Specifically we define 
“high” prior beliefs as believing that the severe adverse events are at least “somewhat 
likely” with vaccination, the median belief. We report the results in Table
\ref{tab:hte_prior}. People without high priors believe side 
effects are lower, and are more likely to intend to vaccinate. We estimate small and 
insignificant  belief updating for this group. For people with high priors, however, we 
estimate much  larger belief updating. The evidence on heterogeneous effects for intentions 
is ambiguous: the point estimates are larger in the high prior group for the industry and 
academic arm butnot the personal arm, and the differences across subgroup are not 
statistically significant. 
We also find that people with diffuse priors update more; Appendix Table
\ref{tab:hte_flu_vacc} shows that people with no flu vaccine experience update more than 
people who vaccinated and experienced no reaction, although they update similarly to people 
experiencing mild reactions.

Table \ref{tab:hte_uni} reports heterogeneous effects by perceived reliability of university
research. People who view university reserach as not reliable or somewhat reliable
show little bleief updating and little resposne of their intentions. People who view
university reserach as reliable show much more belief updating and stronger responses
in their intentions, although the intention respsonse remains small in the industry arm
and the belief response remains small in the representative arm.

These pattern are consistent with belief updating as the key driver of our results 
because we would expect larger belief responses when priors are further from the provided i
nformation. Of course, there is more scope for beliefs to come down when they start at 
a higher level. However, finding university research reliable is associated with lower 
baseline personal side effect beliefs, showing that the level of baseline belief is not 
the only factor driving belief updating.

\subsection{Perceived trial quality drives responses}

We now show that perceived trial quality---trustworthiness and relevance---is a key 
moderator of responses to information. People who view the information as high quality 
show substantial belief updating and increased intention to vaccinate. People who view 
it as low quality show little belief updating and a reduction in their intention to vaccinate.
Like our other heterogeneity analyses, this analysis was not pre-specified and so should
be interpreted with caution.

The average participant did not find the described scientific studies especially
trustworthy or relevant for their decision. In the control arm the average
trustworthiness of the trial was 5.7 and its relevance was 5.11 on a 0-10 scale, and information
provision reduced turstworthiness by 0.3-0.4 points. (Appendix Table \ref{tab:trust_relevance}.)
Negative effects on trustworthiness are consistent with our study sample's high baseline 
prior; seeing information far from a prior makes that information source appaer less 
trustworthy. However, some participants found the trial trustworthy and relevant, and they drive the 
belief updating and intention responses. 

To show the key role of perceived trial quality, we 
would like to estimate heterogeneous effects
by reported trustworthiness and relevance of the trial. However, we assessed trustworthiness
and relevance after providing information, and our information provision affected trial
quality perceptions (Appendix Table \ref{tab:trust_relevance}). To avoid conditioning on an 
intermediate outcome, we do not examine heterogeneous effects by trial quality directly.
Instead we take two steps. 

First, we combine our trust and relevance measures into a single index of perceived
quality. To do so we take the first principal component of both measures, say
$\widetilde{quality}_i$ for an index of perceived quality. Second, to purge 
$\widetilde{quality}_i$ of endogeneity, we work with a prediction given pre-randomization
varaibles only, $\widehat{quality}_i$, obtained as the lasso using all pre-specified
control variables as well as indicators for use of several information sources and
perceived reliability of those sources. We select the lasso penalty parameter using 
10-fold cross-validation. 

We estimate heterogeneous effects by perceived trial quality with the following model:
\begin{equation} \label{eq:hte_quality}
  y_i = \beta_0 + \left(
    \sum_{\text{arm } a} \beta_a 1(arm_i=a) + \gamma_a 1(arm_i=a)\cdot \widehat{quality}_i
    \right)
  + \eta \widehat{quality}_i X_i\theta + \epsilon_i,
\end{equation}
where $X_i$ are our pre-specified controls. Here the effect of arm $a$ for a person
with quality perception $\widehat{quality}$ is $\beta_a + \gamma_a \widehat{quality}$. By
construction, the average effect is $\beta_a$.

We plot the implied effects, by arm, outcome, and quality percetpion, in Figure 
\ref{fig:hte_quality}, 
along with their 95 percent confidence intervals. (Appendix Table \ref{tab:hte_pca} reports
the estimated coefficients and standard errors.) The figure shows substantial
heterogeneity in effects by perceived trial quality. At high quality (a value of 1, 
the 96th percentile), the belief effects are 5-10 percentage points, and the intention effects 
are also large. At low perceived trial quality (a value of -1, the 8th percentile), the beleif 
effects are near zero and the intention effects are negative. This backlash effect
at low perceived quality is consistent with prior work such as \citet{nyhan2014effective},
who also find that providing scientific information (that pediatric vaccines do not cause
autism) reduces vaccine intentions among the most skeptical. The backlash is greatest 
in the industry arm, and this is why the industry arm produces a small average effect on 
intentions. 


